\documentclass[11pt]{article}

\usepackage[letterpaper,margin=1in]{geometry}
\usepackage{amsmath,amssymb}
\usepackage{booktabs}
\usepackage{longtable}
\usepackage{array}
\usepackage{tabularx}
\usepackage{microtype}
\usepackage{placeins}
\usepackage[hidelinks]{hyperref}
\usepackage{xurl}

\newcolumntype{Y}{>{\raggedright\arraybackslash}X}
\newcommand{\model}[1]{\texttt{#1}}
\newcommand{\inprogress}{\textbf{In progress}}

\title{Med-Self-Preference}
\author{
  \textbf{Zara Ansari}\textsuperscript{1,2}\quad
  \textbf{Shlok Natarajan}\textsuperscript{1}\quad
  \textbf{Author Two}\textsuperscript{1}\\
  \textbf{Author Three}\textsuperscript{1}\quad
  \textbf{Aaron Fanous}\textsuperscript{1}\quad
  \textbf{Roxana Daneshjou}\textsuperscript{1,3}\\[0.8em]
  \small \textsuperscript{1}Department of Biomedical Data Science,
  \textsuperscript{3}Department of Dermatology,\\
  \small Stanford University, 450 Jane Stanford Way, Stanford, CA 94305, USA.\\
  \small \textsuperscript{2}Department of Computer Science, Harvey Mudd College,\\
  \small 301 Platt Boulevard, Claremont, CA 91711, USA.\\[0.8em]
  \small Emails: zansari6@stanford.edu, insert email, author2@stanford.edu,\\
  \small author3@stanford.edu, author4@stanford.edu, roxanad@stanford.edu
}
\date{August 2026}

\begin{document}
\maketitle

\begin{abstract}
Benchmarking medical AI has required increasingly larger datasets, making
human evaluation untenable. The common approach now is to test AI models on a
set of tasks and to use large language models (LLMs) as automated judges to
score and compare model performance. However, LLM judges may show
self-preference bias, where a model favors responses produced by its own model.
This has serious implications in medicine since LLM-as-a-judge is readily used
and biased evaluations could affect which models are considered safe.

This study examines LLM-as-a-judge self-preference in two medical settings. We
evaluate eight generators and judges on all 620 Real-POCQi questions and six
clinician models and judges on 200 MedSP1000 standardized-patient scenarios at
two, four, six, and eight visible turns. Model identities are concealed from
the judge in the primary experiments and response order is randomized to
control for positional bias.

In the completed Real-POCQi experiment, judges ranked their own exact
model's answer 1.219 positions higher on average than other judges ranked that
same answer. The effect remained after presentation-position adjustment,
within-judge score calibration, and removal of the explicit rubric. In
MedSP1000, the corresponding effects were 0.518, 0.457, 0.502, and 0.613
positions at two, four, six, and eight turns. Self-preference varied
substantially across judges and did not increase monotonically with every
additional pair of turns. Clinical specialty did not explain significant
heterogeneity in the Real-POCQi effect. These findings show that raw win rate,
score-level preference, and matched rank preference capture different aspects
of self-preference and should be considered together when evaluating LLM
judges in medicine.
\end{abstract}

\section{Introduction}

The complexity of medical practice has called for larger and more complex
benchmarking for comparing the performance of AI models. Human review of
thousands of scenarios tested across multiple AI models is untenable due to
cost and time \cite{zheng2023}. A recent review of medical AI benchmarking
found that large language models (LLMs) are increasingly used as judges to
evaluate the outputs of other LLMs, and in some cases their own outputs
\cite{li2026}. This same study found that OpenAI's models are currently the most
popular judges for medical AI benchmarking studies \cite{li2026}.

Researchers claim using an LLM as a judge provides a scalable way to evaluate
responses based on certain criteria such as accuracy, completeness, safety, and
clarity. The stakes in medical AI evaluation are high, as biased judges could
incorrectly favor unsafe models and influence clinical deployment decisions to
the detriment of patients. Furthermore, evaluating medical responses is
especially complicated since an answer may sound confidently correct, but
still contain incomplete medical information. As a result, evaluating LLMs
used in the medical context is a multifaceted process where faithfulness,
completeness, safety, clarity, and conciseness should be considered.

Prior work in general generative AI tasks has shown that LLM judges may exhibit
biases such as positional and verbosity bias. Wataoka et al. examined how LLM
judges may favor responses that reflect policies or writing styles familiar to
them \cite{wataoka2025}. One study showed that LLM judges may favor
AI-generated clinical responses over human-produced ones. Alvarez-Arenas et
al. found that this preference seemed to depend more on how an answer was
written, especially how long it was and the language used, than on whether the
model could actually tell if the answer came from a human or an AI
\cite{alvarez2026}. Panickssery et al. found that LLM evaluators can recognize
and favor their own generations on general tasks \cite{panickssery2024}.
Pombal et al. found that self-preference can persist even with fully objective
rubrics \cite{pombal2026}.

The literature for self-preference in medical contexts is far more sparse. To
address these gaps, we evaluate self-preference across multiple model families
in both single-turn generalist-specialist question answering and multi-turn
patient-physician conversations. In order to address the fact that a model may
choose its own response because its answer is genuinely better, we hold the
answer fixed and compare how its own model ranks it with how every outside
judge ranks that exact answer. We also evaluate calibration-aware rubric scores
and a no-rubric ranking condition. Our findings provide important
considerations for future medical AI benchmarking studies.

\section{Methods}

\subsection{Generator and judge models}

The study uses two models from each of four families, shown in
Table~\ref{tab:cohort}. These tier labels describe their role in this study and
are not intended as a universal performance ranking.

\begin{table}[htbp]
\centering
\caption{Generator and judge cohort.}
\label{tab:cohort}
\begin{tabular}{lll}
\toprule
Family & Higher-capability model & Workhorse model \\
\midrule
OpenAI & \model{gpt-5.6-sol} & \model{gpt-5.6-terra} \\
Anthropic & \model{claude-opus-5} & \model{claude-sonnet-5} \\
Google & \model{gemini-3.1-pro-preview} & \model{gemini-3.7-flash} \\
Qwen & \model{Qwen/Qwen3.5-122B-A10B-FP8} & \model{Qwen/Qwen3.8-27B-FP8} \\
\bottomrule
\end{tabular}
\end{table}
\FloatBarrier

\subsection{Single-turn pipeline}

The single-turn pipeline uses all 620 questions from the frozen
Real-POCQi source revision. All eight models generated one answer per question,
for 4,960 logical generations. In the primary combined condition, all eight
models judged every question. Each judge received the question and eight
blinded candidate answers in randomized order. It assigned a score from 0 to 5
on accuracy, clinical utility, source quality, verifiability, and completeness,
and then independently ranked all eight answers from best to worst.

A no-rubric direct-ranking sensitivity used the same eight judges on a
deterministic sample of 100 questions. Judges were asked only to rank the eight
blinded candidates by overall quality. An identity-revealed follow-up uses a
seeded 200-question sample and the six API judges; this follow-up was still in
progress at the analysis cutoff.

\subsection{Multi-turn pipeline}

The multi-turn experiment uses the role-separated MedSP1000 dataset. A
deterministic 200-scenario cohort was selected from source
revision \path{55e3e55efd08c73baab912ba0c5b42637114fbc8}. Each scenario has a
private standardized-patient actor packet and a separate clinician-visible
initialization. The patient model receives only the actor packet and visible
conversation; the clinician receives only its initialization and visible
conversation. Evaluator and environment-controller materials are excluded.

The patient simulator is fixed to
\model{mistralai/Mistral-Small-3.1-24B-Instruct-2503}. Six API clinician models
have complete production trajectories: Sol, Terra, Opus, Sonnet, Gemini Pro,
and Gemini Flash. Each trajectory contains four clinician-patient exchanges.
The same trajectories are judged after two, four, six, and eight visible turns.
At each length, all six judges score and rank all six blinded candidates,
yielding 1,200 judgments per length. The two planned Qwen clinician and judge
extensions have not yet been run for MedSP1000.

\section{Analysis of Data}
\subsection{Matched rank self-preference}

For question $q$ and exact model $j$, let $r_{kqj}$ be the rank that judge $k$
assigns to the answer generated by model $j$, where lower ranks are better. The
matched self-preference contrast is
\begin{equation}
D_{jq}=r_{jqj}-\frac{1}{J-1}\sum_{k\ne j}r_{kqj}.
\end{equation}
Thus, $D_{jq}<0$ means that judge $j$ ranks its own answer more favorably than
outside judges rank the same answer to the same question. Matching on the exact
answer removes general generator quality as an explanation.

Candidate order has a measurable association with rank. Ranks are adjusted
with an additive model containing judge-by-generator fixed effects and
categorical presentation-position effects. For pooled intervals, the
correlated judge effects are first averaged within question and uncertainty is
then calculated across questions. Calibration-aware rubric sensitivities
standardize each five-axis rubric sum within every judge-question candidate
list before applying the same matched comparison. Judge-specific families of
tests use Holm correction.

Because the four MedSP1000 conditions are truncations of the same 200
trajectories, transcript-length comparisons are repeated measures. The joint
length test uses the 4-minus-2, 6-minus-2, and 8-minus-2 question-level
contrasts. Individual length contrasts are paired within question.

\section{Results}

\subsection{Real-POCQi generation and aggregate quality}

All 4,960 intended answer cells are present. Claude Opus 5 was materially more
verbose than the other generators, with a median of 2,927 output tokens,
compared with 1,354 for Sonnet and 656--988 for the remaining six models.
Manual review of three matched cases found a substantive venous-reflux error,
an internal HSV-trial contradiction, and omissions or unsupported claims
across models in an oncology question. The generation set is therefore a valid
evaluation input, but not a clinical gold standard.

In the completed combined condition, Opus ranked first overall by a wide
margin. The Qwen generations ranked last, although the Qwen judges preserved
the broad aggregate order.

\begin{table}[htbp]
\centering
\small
\caption{Real-POCQi aggregate generator performance across eight
judges. Lower mean rank is better.}
\label{tab:real-performance}
\begin{tabular}{lrrr}
\toprule
Generator & Mean rank & Ranked first & Mean rubric sum \\
\midrule
Claude Opus 5 & 1.887 & 69.3\% & 23.110 \\
GPT-5.6 Terra & 3.371 & 17.5\% & 21.652 \\
Gemini 3.7 Flash & 3.919 & 1.3\% & 21.079 \\
GPT-5.6 Sol & 4.301 & 8.3\% & 20.838 \\
Claude Sonnet 5 & 4.389 & 1.5\% & 20.632 \\
Gemini 3.1 Pro & 4.620 & 1.3\% & 20.414 \\
Qwen 3.5 122B & 6.524 & 0.4\% & 17.524 \\
Qwen 3.8 27B & 6.989 & 0.3\% & 16.406 \\
\bottomrule
\end{tabular}
\end{table}
\FloatBarrier

\subsection{Real-POCQi self-preference}

Every judge's matched estimate is negative. Every judge-specific interval
excludes zero after Holm correction. The pooled position-adjusted effect is
$-1.219$ rank positions, with a 95\% CI of $[-1.261,-1.177]$. The unadjusted
estimate is also $-1.219$, indicating that position imbalance does not explain
the matched effect.

\begin{table}[htbp]
\centering
\small
\caption{Position-adjusted matched self-preference in Real-POCQi. Negative
values indicate that the judge ranked its own answer more favorably.}
\label{tab:real-self}
\begin{tabular}{lrr}
\toprule
Judge & Rank effect & 95\% CI \\
\midrule
GPT-5.6 Sol & -2.939 & $[-3.041,-2.836]$ \\
GPT-5.6 Terra & -2.224 & $[-2.325,-2.123]$ \\
Gemini 3.7 Flash & -1.236 & $[-1.317,-1.155]$ \\
Claude Opus 5 & -0.969 & $[-1.019,-0.919]$ \\
Qwen 3.5 122B & -0.730 & $[-0.852,-0.608]$ \\
Gemini 3.1 Pro & -0.683 & $[-0.790,-0.575]$ \\
Claude Sonnet 5 & -0.647 & $[-0.763,-0.530]$ \\
Qwen 3.8 27B & -0.324 & $[-0.419,-0.228]$ \\
\midrule
Pooled within question & -1.219 & $[-1.261,-1.177]$ \\
\bottomrule
\end{tabular}
\end{table}
\FloatBarrier

The Qwen results show why raw first-place rate is insufficient. Qwen answers
remain weak in absolute terms, but each Qwen judge moves its own answer upward
relative to how the outside judges rank that same answer. A
calibration-aware rubric sensitivity gives a pooled effect of $-0.440$
within-list standard deviation units (95\% CI $[-0.460,-0.420]$), with all
eight judge-specific effects negative.

\subsection{No-rubric direct-ranking sensitivity}

On the deterministic 100-question sample, the pooled matched effect is
$-1.263$ positions (95\% CI $[-1.361,-1.164]$). Seven judge-specific intervals
exclude zero; Qwen 27B is the exception at $-0.155$ positions (95\% CI
$[-0.346,0.037]$). Direct and rubric-condition rankings agree on 91.5\% of
candidate pairs, select the same first-place answer in 87.1\% of cells, and
produce the identical complete order in 16.1\%. Removing the explicit rubric
therefore changes fine ordering but does not remove the main self-preference
pattern.

\subsection{Clinical specialty analysis}

All adequately sized specialties have negative point estimates. Among
specialties with at least ten questions, the pooled effects range from
$-1.543$ positions in Family Medicine to $-0.908$ in Radiology. However, a
label-permutation test preserving the 30 observed specialty sizes finds that
specialty explains only 5.7\% of question-level variation and does not show
significant global heterogeneity ($p=0.190$). No judge-specific specialty test
crosses 0.05. The correct interpretation is not that every specialty has an
identical true effect, but that the current data do not provide convincing
evidence of a specialty interaction.

\subsection{MedSP1000 results}

All four transcript-length conditions are complete at 1,200 judgments each.
Matched self-preference is present at every length. The calibration-aware
rubric sensitivity points in the same direction.

\begin{table}[htbp]
\centering
\small
\caption{Pooled MedSP1000 matched self-preference by visible transcript length.}
\label{tab:med-pooled}
\begin{tabular}{rrrr}
\toprule
Visible turns & Rank effect & 95\% CI & Normalized-rubric effect \\
\midrule
2 & -0.518 & $[-0.585,-0.452]$ & -0.309 \\
4 & -0.457 & $[-0.527,-0.388]$ & -0.305 \\
6 & -0.502 & $[-0.574,-0.430]$ & -0.356 \\
8 & -0.613 & $[-0.683,-0.543]$ & -0.405 \\
\bottomrule
\end{tabular}
\end{table}
\FloatBarrier

Judge-specific patterns vary over length and do not share a simple trajectory.
Opus becomes more self-preferring with length, Gemini Pro is strongest at two
turns, and Sol has a marked increase only at eight turns.

\begin{table}[htbp]
\centering
\small
\caption{Judge-specific MedSP1000 matched rank effects.}
\label{tab:med-judge}
\begin{tabular}{lrrrr}
\toprule
Judge & 2 turns & 4 turns & 6 turns & 8 turns \\
\midrule
Claude Opus 5 & -0.153 & -0.218 & -0.423 & -0.445 \\
Claude Sonnet 5 & -0.233 & -0.378 & -0.549 & -0.472 \\
Gemini 3.1 Pro & -0.791 & -0.554 & -0.436 & -0.545 \\
Gemini 3.7 Flash & -0.572 & -0.600 & -0.428 & -0.645 \\
GPT-5.6 Sol & -0.536 & -0.487 & -0.539 & -0.938 \\
GPT-5.6 Terra & -0.826 & -0.508 & -0.637 & -0.632 \\
\bottomrule
\end{tabular}
\end{table}
\FloatBarrier

A repeated-measures test rejects equality across the four lengths
($F(3,197)=3.972$, $p=0.0089$). The 4-minus-2 and 6-minus-2 changes are small
and do not exclude zero. The 8-minus-2 change is $-0.094$ positions (95\% CI
$[-0.187,-0.002]$, nominal $p=0.045$), but does not survive Holm correction
across the three baseline contrasts. The adjacent 8-minus-6 change is
$-0.111$ positions (95\% CI $[-0.189,-0.033]$, nominal $p=0.0054$). The
defensible conclusion is that self-preference varies with visible length and is
strongest at eight turns in the pooled data, not that it rises monotonically
from two through eight turns.

\subsection{Experiments still in progress or not completed}

The identity-revealed Real-POCQi experiment was actively running at the fixed
August 24 snapshot. It had 789 of 1,200 planned latest successful judgments,
with 128 of 200 questions complete for all six API judges. Interim results are
not treated as final in the main conclusions. The MedSP1000 Qwen expansion,
production direct-ranking condition, and production rubric-sum-only condition
have not been run. Physician validation of the judging conditions also remains
incomplete.

\section{Discussion}

As medical benchmarking and monitoring relies more on LLMs-as-a-judge, it has
become increasingly important to understand how LLMs exhibit
self-preference. Our study focuses on healthcare, where benchmarking results
can impact which models are selected for implementation. We examine
self-preference across different model families, separating scoring generosity
from direct preferences and from a matched comparison in which the exact answer
is held fixed.

The experiments show that self-preference can look different depending on how
it is measured. Every Real-POCQi judge
ranks its own answer more favorably than outside judges rank the same answer,
including models whose answers rarely rank first. This effect remains after
position adjustment, within-judge score calibration, and removal of the
explicit rubric. It is safer to describe this as same-model affinity rather
than conscious self-recognition. Generator identities were hidden, so the
effect may arise from shared style, reasoning conventions, favored evidence,
or family-specific quality criteria.

The multi-turn findings qualify the role of conversation length.
Self-preference is present at every tested MedSP1000
length and is strongest at eight turns in the pooled data, but it does not
steadily increase from two to four to six to eight turns. Different judges show
different trajectories. Single-turn benchmarks may still miss behavior that
appears in longer conversations, but the length effect should not be described
as universally monotonic.

Clinical specialty does not explain reliable heterogeneity in the current
Real-POCQi sample. The magnitude is driven much more by which model is judging.
This suggests that medical benchmarking studies should report results by judge
and use multiple judge families rather than relying on one pooled evaluator.

\subsection{Limitations}

There are limitations to our study. The experiments have one saved generation
per model-question cell and one
successful judging realization per logical cell. The analysis therefore cannot
separate stable model behavior from generation-specific or run-to-run
idiosyncrasy.

Rankings are mutually dependent within each candidate list. Question-level
intervals address clustering for the pooled estimand, but remain large-sample
approximations. Exact-model and broader family affinity are also difficult to
separate completely. Response length and stylistic differences may mediate the
effect. Opus is especially verbose in Real-POCQi, and clinician output length
also differs in MedSP1000.

The specialty labels are broad and highly imbalanced, ranging from three to 55
questions. The absence of detected heterogeneity is not proof that narrower
clinical domains have identical effects. The identity-revealed experiment is
in progress, the full MedSP1000 Qwen extension has not been run, and the
automated judges have not yet been benchmarked against completed blinded
physician judgments. Finally, the analysis was exploratory and was
not preregistered. Effect sizes and sensitivity checks should be emphasized
over isolated p-values.

\section{Conclusion}

Overall, our findings show that self-preference is present in medical LLM
judging, but it does not appear in the same way across all models or metrics.
The experiments show that LLM judges rank answers from their own exact
model more favorably than other judges rank the same answers. This pattern is
robust across two medical settings, position adjustment, calibrated rubric
scores, and a no-rubric sensitivity analysis. In multi-turn conversations it
is present at every tested length and strongest at eight turns in pooled data,
without a monotonic increase from two to eight turns. These results highlight
the importance of using multiple judge families, multiple evaluation measures,
and matched analyses when comparing medical LLMs.

\clearpage
\appendix

\section{Complete experiment status}

\begin{longtable}{>{\raggedright\arraybackslash}p{0.21\textwidth}
                  >{\raggedright\arraybackslash}p{0.42\textwidth}
                  >{\raggedright\arraybackslash}p{0.25\textwidth}}
\caption{Status of experiments included in the project at the fixed analysis
snapshot.}\label{tab:status}\\
\toprule
Phase & Experiment & Status \\
\midrule
\endfirsthead
\toprule
Phase & Experiment & Status \\
\midrule
\endhead
Current & Eight-model Real-POCQi generation, 4,960 answers & Complete \\
Current & Blinded Real-POCQi combined judging, 4,960 judgments & Complete \\
Current & Real-POCQi no-rubric direct ranking, 800 judgments & Complete \\
Current & Real-POCQi specialty analysis & Complete \\
Current & Identity-revealed Real-POCQi, 1,200 planned judgments & \inprogress; 789 successes at snapshot \\
Current & Six-model MedSP1000 generation, 1,200 trajectories & Complete \\
Current & MedSP1000 judging at 2, 4, 6, and 8 turns, 1,200 each & Complete \\
Planned & MedSP1000 Qwen clinician and judge extension & Not started \\
Planned & MedSP1000 production direct ranking and rubric-sum-only conditions & Pilot only \\
Validation & Blinded physician judging benchmark & Not completed \\
\bottomrule
\end{longtable}

\section{Real-POCQi specialty estimates}

\begin{longtable}{lrrr}
\caption{Position-adjusted matched self-preference by clinical specialty.
Negative values indicate self-preference.}\label{tab:specialty}\\
\toprule
Specialty & Questions & Rank effect & Rubric-rank sensitivity \\
\midrule
\endfirsthead
\toprule
Specialty & Questions & Rank effect & Rubric-rank sensitivity \\
\midrule
\endhead
Endocrinology & 55 & -1.359 & -1.320 \\
Cardiology & 54 & -1.189 & -1.064 \\
Oncology / Hematology & 49 & -1.052 & -1.029 \\
OB / GYN & 44 & -1.283 & -1.216 \\
Dermatology & 40 & -1.097 & -1.113 \\
Rheumatology & 39 & -1.234 & -1.166 \\
Neurology & 32 & -1.231 & -1.158 \\
Infectious Disease & 30 & -1.322 & -1.348 \\
Psychiatry & 25 & -1.359 & -1.244 \\
Surgery & 25 & -1.371 & -1.225 \\
Nephrology & 20 & -1.172 & -1.048 \\
Allergy \& Immunology & 19 & -1.249 & -1.128 \\
Pediatrics & 19 & -1.188 & -1.193 \\
Internal Medicine & 17 & -1.256 & -1.235 \\
Critical Care & 14 & -1.221 & -1.197 \\
Gastroenterology & 14 & -1.141 & -1.101 \\
Pulmonology & 14 & -1.197 & -1.065 \\
Transplantation & 14 & -1.140 & -1.044 \\
Anesthesiology & 11 & -1.297 & -1.091 \\
Family Medicine & 11 & -1.543 & -1.308 \\
Orthopedics & 11 & -1.215 & -1.265 \\
Primary Care & 11 & -1.397 & -1.258 \\
Emergency Medicine & 10 & -1.243 & -1.232 \\
Radiology & 10 & -0.908 & -1.023 \\
Urology & 8 & -1.119 & -0.897 \\
Hospital-Based Medicine & 7 & -0.862 & -0.684 \\
Surgical Oncology & 6 & -1.063 & -0.913 \\
Genetics & 4 & -0.684 & -0.645 \\
Otolaryngology & 4 & -0.933 & -0.799 \\
Ophthalmology & 3 & -1.056 & -0.997 \\
\bottomrule
\end{longtable}

\begin{thebibliography}{9}

\bibitem{zheng2023}
Zheng, L., Chiang, W.-L., Sheng, Y., et al. (2023).
Judging LLM-as-a-Judge with MT-Bench and Chatbot Arena.
\emph{Advances in Neural Information Processing Systems}, 36.

\bibitem{li2026}
Li, L., Li, D., Chen, C., et al. (2026).
LLM-as-a-Judge in Healthcare: A Scoping Analysis of Applications, Methods, and
Human Alignment. \emph{arXiv preprint arXiv:2605.25273}.

\bibitem{wataoka2025}
Wataoka, K., Takahashi, T., and Ri, R. (2025).
Self-Preference Bias in LLM-as-a-Judge.
\emph{arXiv preprint arXiv:2410.21819}.

\bibitem{alvarez2026}
Alvarez-Arenas, J. I., Jimenez-Carretero, D., Mananes, D., and
Sanchez-Cabo, F. (2026).
The Unreliable Judges: Assessing Reproducibility and Self-Preference Bias of
LLMs as Free-Text Evaluators. \emph{medRxiv}.

\bibitem{pombal2026}
Pombal, J., Rei, R., and Martins, A. F. T. (2026).
Self-Preference Bias in Rubric-Based Evaluation of Large Language Models.
\emph{arXiv preprint arXiv:2604.06996}.

\bibitem{panickssery2024}
Panickssery, A., Bowman, S. R., and Feng, S. (2024).
LLM Evaluators Recognize and Favor Their Own Generations.
\emph{Advances in Neural Information Processing Systems}, 37.

\end{thebibliography}

\end{document}
